\begin{textsample}{POS Dim 2 – 2022\_c4 – Score 22.00 – 2022\_c4/enade\_2022\_tec\_gest\_com\_7f0e07f670.txt}  \label{ex:f2_pos_027}
CAPÍTULO 5 PERCEPÇÃO DA PROVA Nas análises feitas neste capítulo, constam as percepções dos concluintes da área de Tecnologia em Gestão Comercial sobre a prova aplicada no Enade/2022.
Essas percepções foram mensuradas por meio de nove questões que avaliaram desde o grau de dificuldade da prova até o tempo gasto para concluí -la.
As percepções sobre a prova foram relacionadas ao desempenho dos estudantes, à grande região de funcionamento do curso, à categoria administrativa e ao tipo de organização acadêmica da IES.
O \textbf{questionário} de percepção da prova encontra -se ao final do Anexo VIII, no qual está a \textbf{reprodução} do exame.
Na apresentação dos dados relativos às nove questões sobre as percepções a respeito da prova, o desempenho dos estudantes foi classificado em quatro quartos.
Para tanto, esse desempenho foi \textbf{ordenado} de forma ascendente.
O percentil 25, P25, também conhecido como primeiro quartil, é a nota de desempenho que deixa um quarto ( 25 \% ) dos valores observados abaixo e três quartos acima.
O quarto inferior de desempenho é composto pelas notas abaixo do primeiro quartil.
Já o percentil 75, P75, também conhecido como terceiro quartil, é o valor para o qual há três quartos ( 75 \% ) dos dados abaixo e um quarto acima dele.
O quarto superior de desempenho é composto pelas notas iguais ou acima do terceiro quartil.
O percentil 50, P50, também conhecido como mediana, é o valor que divide as notas em dois conjuntos de igual tamanho.
O 2.º quarto inclui os valores entre o primeiro quartil ( P25 ) e a mediana.
O terceiro quarto contém os valores entre a mediana ( P50 ) e o terceiro quartil ( P75 ).
Vale ressaltar que percentis, quartis e medianas são pontos que, não obrigatoriamente, pertencem ao conjunto original de dados, ao passo que os quartos são \textbf{subconjuntos} dos dados originais.
Na Figura 5.1, apresenta -se uma ilustração dos quatro quartos descritos.
A seguir, serão apresentados gráficos com resultados selecionados, relativos às nove questões avaliadas por grupos de estudantes.
Nas barras dos gráficos, apresenta -se o percentual de estudantes que assinalaram uma das opções ou a soma das porcentagens daqueles que assinalaram duas ( ou três ) delas.
Por exemplo, para as questões 1 e 2, nos gráficos, é apresentada a porcentagem total de participantes que assinalaram as opções ( D ) “ Difícil ” e ( E ) “ Muito difícil ”.
Em cada barra, foram assinalados, também, os extremos do intervalo de confiança de 95 \% como linhas verticais unidas por uma linha horizontal na forma da letra H maiúscula.
O \textbf{estimador} de um \textbf{parâmetro} com um certo nível de confiança ( por exemplo, 95,0 \% ) deve conter o \textbf{parâmetro} no intervalo de confiança em 95 \% das vezes.
Na comparação entre os \textbf{estimadores} dos \textbf{parâmetros} de duas classes de uma dada categoria ( por exemplo : Norte e Nordeste nas grandes regiões, de primeiro e último quarto dentro de desempenho, de IES públicas e privadas, ou de universidades e faculdades ), associados aos seus respectivos intervalos de confiança, diz -se que não há diferença estatisticamente significativa entre os \textbf{parâmetros} das duas categorias quando há \textbf{interseção} entre os intervalos de confiança, e que há diferença, se os intervalos de confiança são \textbf{disjuntos} ( para maiores informações \textbf{vide} Glossário ).
Nas Tabelas do Anexo V, são apresentados os valores absolutos e a distribuição percentual9 das alternativas válidas das nove questões, segundo o mesmo \textbf{recorte} de desempenho dos estudantes e da grande região de funcionamento do curso, de categoria administrativa e do tipo de organização acadêmica da IES. 9 \textbf{Reitera} -se que uma das convenções para tabelas \textbf{numéricas} ( Subseção 2.2.3 do Capítulo 2 ) refere -se à possibilidade de, por questão de \textbf{arredondamento}, a soma das partes não resultar em 100 \%. 5.1 GRAU DE DIFICULDADE DA PROVA Nesta seção são mostradas as percepções dos estudantes quanto ao grau de dificuldade da prova Enade/2022, e que foram analisadas considerando sua vinculação ao Componente de Formação Geral e ao Componente Específico. 5.1.1 COMPONENTE DE FORMAÇÃO GERAL O Gráfico 5.1 apresenta as respostas obtidas para a Questão 1 : “ Qual o grau de dificuldade desta prova na parte de Formação Geral? ”.
De acordo com os dados do Gráfico 5.1, a região Sudeste apresentou o maior percentual de estudantes que consideraram a prova “ Difícil ” ou “ Muito difícil ”, com 33 \%, seguida da região Centro-Oeste com 32,4 \%.
A região de menor proporção foi a Sul.
Na região Nordeste, a média foi de 30,7 \%, na Norte 29 \% e na região Sul 28,5 \%.
O maior intervalo foi observado na região Norte.
O Gráfico 5.2 traz mais resultados sobre a questão tratada no gráfico anterior, porém por quartos de desempenho.
Os dados do Gráfico 5.2 mostram que o percentual de estudantes que consideraram o Componente de Formação Geral da prova “ Difícil ” ou “ Muito difícil ” foi decrescente em função dos quartos de desempenho : 38,2 \%, no 1º quarto, e 24 \%, no 4º quarto.
Nos quartos de desempenho intermediários, a proporção de estudantes que consideraram a prova “ Difícil ” ou “ Muito difícil ” foi de 36,6 \%, no 2º quarto, e de 29,4 \%, no 3º quarto.
O Gráfico 5.3 traz outros resultados também em relação à mesma questão retratada no gráfico anterior, porém por organização acadêmica das IES.
Como mostram os dados do Gráfico 5.3, o percentual de estudantes que consideraram o Componente de Formação Geral da prova “ Difícil ” ou “ Muito difícil ” foi maior entre os estudantes de CEFET/IF, alcançando 34,5 \%, enquanto entre os estudantes das faculdades, a organização acadêmica com a menor incidência, o percentual alcançou 28,7 \%.
O Gráfico 5.4 apresenta os dados por categoria administrativa das IES.
Os dados do Gráfico 5.4 mostram que o percentual de estudantes que consideraram o Componente de Formação Geral da prova “ Difícil ” ou “ Muito difícil ” foi maior entre os estudantes de cursos de IES privadas ( 32,2 ). 5.1.2 COMPONENTE ESPECÍFICO O Gráfico 5.5 apresenta as respostas obtidas para a Questão 2 : “ Qual o grau de dificuldade desta prova na parte de Componente Específico? ”.
De acordo com os dados do Gráfico 5.5, a região Norte apresentou o maior percentual de estudantes que consideraram o Componente Específico da prova “ Difícil ” ou “ Muito difícil ”, com 36,8 \%, seguida da região Centro-Oeste com 36,1 \%.
Na região Nordeste, a média foi de 33,1 \%, na Sudeste 34,9 \% e na região Sul 29,3 \%.
O Gráfico 5.6 traz mais resultados sobre a questão tratada no gráfico anterior, porém por quartos de desempenho.
Os dados do Gráfico 5.6 mostram que o percentual de estudantes que consideraram o Componente Específico da prova “ Difícil ” ou “ Muito difícil ” foi decrescente em função dos quartos de desempenho : 38,7 \%, no 1º quarto, e 29,3 \%, no 4º quarto, o grupo de melhor desempenho na prova.
Nos quartos de desempenho intermediários, a proporção de estudantes que consideraram a prova “ Difícil ” ou “ Muito difícil ” foi de 36,8 \%, no 2º quarto, e de 32,5 \% no 3º quarto.
O Gráfico 5.7 traz outros resultados também em relação à mesma questão retratada no gráfico anterior, porém por organização acadêmica das IES.
Como mostram os dados do Gráfico 5.7, o percentual de estudantes que consideraram o Componente de Formação Geral da prova “ Difícil ” ou “ Muito difícil ” foi maior entre os estudantes de CEFET/IF, alcançando 41,2 \%, enquanto entre os estudantes das faculdades, a organização acadêmica com a menor incidência, o percentual alcançou 30,1 \%.
O Gráfico 5.8 apresenta os dados por categoria administrativa das IES.
Os dados do Gráfico 5.8 mostram que o percentual de estudantes que consideraram o Componente Específico da prova “ Difícil ” ou “ Muito difícil ” foi ligeiramente maior entre os estudantes de cursos de IES privadas ( 34,3 \% ).
Enquanto a proporção registrada em IES públicas foi de 32,9 \%. 5.2 EXTENSÃO DA PROVA EM RELAÇÃO AO TEMPO TOTAL O Gráfico 5.9 apresenta as respostas obtidas para a Questão 3 : “ Considerando a extensão da prova, em relação ao tempo total, você considera que a prova foi longa ou muito longa? ”.
Em relação aos dados do Gráfico 5.9, entre as grandes regiões, os que avaliaram a prova como longa ou muito longa, em relação ao tempo total destinado à sua resolução, atingiram uma porcentagem que variou de 31,7 \%, na região Nordeste, até 45 \%, na região Sul.
O Gráfico 5.10 traz mais resultados sobre a questão tratada no gráfico anterior, porém por quartos de desempenho.
No Gráfico 5.10, observa -se que a proporção de estudantes que consideraram a prova longa ou muito longa foi menor no 3º quarto de desempenho ( 34,7 \% ), e maior, no 1º quarto ( 40,2 \% ).
O Gráfico 5.11 traz outros resultados também em relação à mesma questão retratada no gráfico anterior, porém por organização acadêmica das IES.
De acordo com os dados do Gráfico 5.11, considerando -se a extensão da prova, o percentual de estudantes que avaliaram a prova como longa ou muito longa, em relação ao tempo total destinado à sua resolução, foi maior entre os estudantes de faculdades, alcançando 42,7 \%, enquanto entre os estudantes de CEFET/IF, a organização acadêmica com a menor incidência, esse índice atingiu 27,7 \%.
Entre os participantes de centros universitários, o índice atingiu 35,7 \%, e entre os de universidades, 36,7 \%.
O Gráfico 5.12 apresenta os dados por categoria administrativa das IES.
Os dados do Gráfico 5.12 mostram que o percentual de estudantes que consideraram a prova longa ou muito longa, em relação ao tempo total destinado à sua resolução, foi maior entre os estudantes de cursos de IES privadas ( 36,9 \% ) do que entre os estudantes de IES públicas ( 33,5 \% ). 5.3 COMPREENSÃO DOS ENUNCIADOS DAS QUESTÕES Também organizados em relação ao Componente de Formação Geral e ao Componente Específico, os dados apresentados nesta seção dizem respeito à forma como os estudantes compreenderam os enunciados da questões da prova, considerando as formulações expressas nesses mesmos enunciados. 5.3.1 COMPONENTE DE FORMAÇÃO GERAL O Gráfico 5.13 apresenta as respostas obtidas para a Questão 4 : “ Os enunciados das questões da prova na parte de Formação Geral estavam claros e objetivos? ”.
Os resultados apresentados no Gráfico 5.13 mostam que, na análise regional, a \textbf{percentagem} de estudantes que avaliaram que todos ou a maioria dos enunciados das questões do Componente de Formação Geral estavam “ claros e objetivos ” variou de 76,2 \%, na região Centro-oeste e 82,7 \%, na região Norte.
O Gráfico 5.14 traz mais resultados sobre a questão tratada no gráfico anterior, porém por quartos de desempenho.
De acordo com os dados do Gráfico 5.14, segundo o desempenho, observa -se que a proporção dos que emitiram essa opinião \textbf{cresce}, conforme o desempenho aumenta.
No quarto superior, a clareza e a objetividade de todos ou da maioria dos enunciados das questões foram percebidas por 87,1 \%.
O Gráfico 5.15 traz outros resultados também em relação à mesma questão retratada no gráfico anterior, porém por organização acadêmica das IES.
Os dados do Gráfico 5.15 mostram que a proporção dos presentes que consideraram que todos ou a maioria dos enunciados das questões do Componente de Formação Geral estavam “ claros e objetivos ” foi maior entre os estudantes de faculdades, chegando a 82,7 \%, enquanto entre CEFET/IF, a organização acadêmica com a menor incidência, esse índice atingiu 76,4 \%.
Entre os estudantes de centros universitários, o índice atingiu 78,8 \%, e entre as universidades, 77,2 \%.
O Gráfico 5.16 apresenta os dados por categoria administrativa das IES.
Os dados do Gráfico 5.16 mostram que o percentual de estudantes que consideraram que todos ou a maioria dos enunciados das questões do Componente de Formação Geral estavam “ claros e objetivos ” foi maior entre os estudantes de cursos de IES públicas ( 82,4 \% ) quando comparado aos estudantes de IES privadas ( 77,7 \% ). 5.3.2 COMPONENTE ESPECÍFICO O Gráfico 5.17 apresenta as respostas obtidas para a Questão 5 : “ Os enunciados das questões da prova na parte de Componente Específico estavam claros e objetivos? ”.
A maioria dos estudantes de todas as grandes regiões brasileiras considerou “ claros e objetivos ” todos ou a maioria dos enunciados das questões do Componente Específico da prova, percentual sempre maior ou igual a 76,3 \%.
O Gráfico 5.18 traz mais resultados sobre a questão tratada no gráfico anterior, porém por quartos de desempenho.
Os dados do Gráfico 5.18 mostram que a proporção de estudantes que consideraram os enunciados das questões “ claros e objetivos ” apresenta uma tendência crescente em relação ao aumento de desempenho : mais elevada no quarto superior ( 88,2 \% ), se comparada ao quarto inferior de desempenho ( 63,8 \% ).
O Gráfico 5.19 traz outros resultados também em relação à mesma questão retratada no gráfico anterior, porém por organização acadêmica das IES.
Os dados do Gráfico 5.19 mostram que a proporção dos presentes que consideraram que todos ou a maioria dos enunciados das questões do Componente Específico estavam “ claros e objetivos ” foi maior entre os estudantes de faculdades, atingindo 84,4 \%.
Entre os estudantes de universidades e de centros universitários, ambas organizações acadêmicas com a menor incidência, essa proporção atingiu 76 \% e 78,6 \%, respectivamente, sendo intermediária a proporção alcançada entre os estudantes de CEFET/IF ( 81 \% ).
O Gráfico 5.20 apresenta os dados por categoria administrativa das IES.
Os dados do Gráfico 5.20 mostram que o percentual de estudantes que consideraram que todos ou a maioria dos enunciados das questões do Componente Específico estavam “ claros e objetivos ” foi maior entre os estudantes de cursos de IES públicas ( 82,6 \% ) que entre os de IES privadas ( 77,2 \% ). 5.4 SUFICIÊNCIA DAS INFORMAÇÕES/INSTRUÇÕES FORNECIDAS O Gráfico 5.21 apresenta as respostas obtidas para a Questão 6 : “ As informações/instruções fornecidas para a resolução das questões foram suficientes para resolvê -las? ”.
De acordo com os dados do Gráfico 5.21 e quanto à distribuição de respondentes pelas grandes regiões, observa -se que a proporção de estudantes que consideraram as informações/instruções fornecidas até excessivas ou suficientes em todas ou na maioria das questões foi sempre superior ou igual a 83,5 \%, chegando a 89,1 \%, na região Sul.
O Gráfico 5.22 traz resultados sobre a questão tratada no gráfico anterior, porém por quartos de desempenho.
De acordo com os dados do Gráfico 5.22, levando -se em conta o desempenho dos participantes, observa -se uma tendência crescente das proporções de participantes que avaliaram as informações/instruções como até excessivas ou suficientes em todas ou na maioria das questões.
O percentual foi mais elevado no quarto superior ( 91,9 \% ).
Já no quarto inferior, a suficiência das informações/instruções declarada como até excessiva, em todas ou na maioria das questões foi percebida por 76,1 \% dos respondentes.
O Gráfico 5.23 traz outros resultados também em relação à mesma questão retratada no gráfico anterior, porém por organização acadêmica das IES.
De acordo com os dados do Gráfico 5.23, considerando -se as informações/instruções fornecidas, a proporção que considerou que essas eram até excessivas ou suficientes em todas ou na maioria das questões foi maior entre os estudantes de faculdades, \textbf{totalizando} 90,9 \%.
Os CEFET/IF foram a organização acadêmica em que houve a menor proporção : 81,6 \%.
O Gráfico 5.24 apresenta os dados por categoria administrativa das IES.
De acordo com os dados do Gráfico 5.24, o percentual de estudantes que consideraram que as informações/instruções fornecidas para a resolução das questões eram até excessivas ou suficientes em todas ou na maioria das questões foi maior entre os estudantes de cursos de IES públicas ( 86,9 \% ), não havendo diferença estatisticamente significativa da proporção registrada entre os de IES privadas ( 85,3 \% ). 5.5 DIFICULDADE ENCONTRADA AO RESPONDER À PROVA Nos Gráficos 5.25 a 5.28, apresentados de forma \textbf{sequenciada} mais adiante, são expostos os percentuais de estudantes que apontaram o desconhecimento do conteúdo como “ dificuldade ao responder à prova ”.
O Gráfico 5.25 apresenta as respostas obtidas sobre o desconhecimento do conteúdo como a principal “ dificuldade ao responder à prova ”.
De acordo com os dados o Gráfico 5.25, os percentuais variaram de 11,4 \%, nas regiões Sul e Nordeste, a 14,5 \%, na região Centro-oeste.
O Gráfico 5.26 traz resultados sobre o ponto tratado no gráfico anterior, porém por quartos de desempenho.
De acordo com os dados do Gráfico 5.26, em relação aos quartos de desempenho, o desconhecimento do conteúdo foi a opção escolhida por 14,9 \% dos estudantes do quarto inferior e por 10,2 \% do quarto superior, caracterizando uma tendência decrescente.
O Gráfico 5.27 traz outros resultados também em relação à mesma questão retratada no gráfico anterior, porém por organização acadêmica das IES.
De acordo com os dados do Gráfico 5.27, na análise por organização acadêmica, os percentuais variaram de 11,5 \%, nos CEFET/IF, a 13,7 \%, nos centros universitários.
O Gráfico 5.28 apresenta os dados por categoria administrativa das IES.
De acordo com os dados do Gráfico 5.28, em relação às categorias administrativas, o Desconhecimento do conteúdo foi a opção escolhida por 12,8 \% dos estudantes, nas IES privadas, percentual ligeiramente maior do que nas IES públicas ( 12,5 \% ). 5.6 CONTEÚDOS DAS QUESTÕES OBJETIVAS DA PROVA O Gráfico 5.29 apresenta as respostas obtidas para a Questão 8 : “ Considerando apenas as questões objetivas da prova, você percebeu que não estudou ainda a maioria desses conteúdos? ”.
Na análise por grande região, a proporção de respondentes que escolheram a opção “ não estudou ainda a maioria desses conteúdos ” foi pequena.
Observa -se que, nas regiões Norte ( 6,5 \% ) e Centro-oeste ( 8,2 \% ), as proporções foram ligeiramente maiores.
Nas regiões Sul e Nordeste, a situação foi \textbf{inversa}, com proporções mais baixas, correspondendo a, respectivamente, 3,2 \% e 4,7 \%.
A região Sudeste ficou com a proporção de 5,3 \%.
O Gráfico 5.30 traz resultados sobre a questão tratada no gráfico anterior, porém por quartos de desempenho.
De acordo com os dados do Gráfico 5.30, considerando -se separadamente as opiniões de estudantes dos quatro quartos de desempenho, observa -se que, no quarto inferior, 10,6 \% ofereceram como resposta que não estudaram ainda a maioria desses conteúdos, sendo 2,4 \% os do quarto superior com a mesma resposta.
O Gráfico 5.31 traz outros resultados também em relação à mesma questão retratada no gráfico anterior, porém por organização acadêmica das IES.
De acordo com os dados do Gráfico 5.31, na análise por organização acadêmica, a proporção de respondentes que escolheram a opção “ não estudou ainda a maioria desses conteúdos ” foi pequena.
Observa -se que as proporções nas faculdades ( 3,6 \% ) e nas universidades ( 5,4 \% ) foram menores.
Nos centros universitários e nos CEFET/IF, a situação foi \textbf{inversa} : a proporção foi mais alta, 5,8 \% e 6,1 \%, respectivamente.
O Gráfico 5.32 apresenta os dados por categoria administrativa das IES.
De acordo com os dados do Gráfico 5.32, em relação às categorias administrativas, observa -se que, nas IES públicas, 4,5 \% ofereceram como resposta que “ não estudou ainda a maioria desses conteúdos ”, sendo 5,5 \% nas IES privadas com a mesma resposta. 5.7 TEMPO GASTO PARA CONCLUIR A PROVA O Gráfico 5.33 apresenta as respostas obtidas para a Questão 9 : “ Qual foi o tempo gasto por você para concluir a prova? ”.
Considerando -se as cinco grandes regiões brasileiras, os percentuais dos que utilizaram entre duas e quatro horas para finalizar a prova foram, nas regiões Norte e Nordeste, respectivamente 75,3 \% e 74,3 \%.
Nas regiões Sul, Sudeste e Centro-Oeste, os percentuais de estudantes que \textbf{dispensaram} entre duas e quatro horas para concluir a prova ficaram muito próximos, 70,4 \%, 71,6 \% e 69,3 \%, respectivamente.
O Gráfico 5.34 traz resultados sobre a questão tratada no gráfico anterior, porém por quartos de desempenho.
De acordo com os dados do Gráfico 5.34, uma vez analisadas as alternativas escolhidas pelos estudantes que se situam nos diferentes quartos de desempenho, observa -se uma tendência crescente : uma proporção maior de participantes no quarto superior declarou ter gasto entre duas e quatro horas para concluir a prova quando comparada com a proporção de participantes nos quartos inferiores.
O Gráfico 5.35 traz outros resultados também em relação à mesma questão retratada no gráfico anterior, porém por organização acadêmica das IES.
De acordo com os dados do Gráfico 5.35, na análise por organização acadêmica, os percentuais dos que utilizaram entre duas e quatro horas para finalizar a prova foram de 74,8 \% nas faculdades ; 71,7 \% nas universidades e 78,2,7 \% nos CEFET/IF.
O menor percentual foi entre os estudantes de centros universitários, com 71 \% dos alunos necessitando de duas a quatro horas para concluir a prova.
O Gráfico 5.36 apresenta os dados por categoria administrativa das IES.
De acordo com os dados do Gráfico 5.36, em relação às categorias administrativas, observa -se que, nas IES públicas, 78 \% declararam ter gasto entre duas e quatro horas para concluir a prova, sendo 71,2 \% nas IES privadas com a mesma resposta.
Os resultados do Enade 2022 apresentados neste capítulo contribuem para o entendimento de que a consideração da percepção dos estudantes concluintes sobre a prova do Enade é importante tanto para a compreensão do seu desempenho como para a reflexão sistemática e constante a respeito de determinadas realidades relacionadas a esse mesmo desempenho, no país.
Por outro lado, como todos os elementos do trabalho \textbf{pedagógico} institucional mostram -se diretamente relacionados com o desempenho dos estudantes, também ganha relevância conhecer e analisar os resultados da percepção dos \textbf{coordenadores} de curso sobre a mesma prova.
As informações expressas por estudantes e \textbf{coordenadores} sobre questões diversas associadas aos dois componentes da prova – Formação Geral e Conhecimento Específico – podem agregar valor em ações instituionais voltadas para o aperfeiçoamento do instrumento \textbf{avaliativo}, considerndo os vários elementos envolvido em sua elaboração.
Neste contexto, é relevante que os resultados das percepções de estudantes e \textbf{coordenadores} sobre a prova do Enade ganhem peso como um elemento a mais nas análises da oferta dos cursos nas modalidades presencial e a distância, em nível nacional.
Espera -se, assim, que, em uma visão formativa de avaliação, o conhecimento e a análise desses resultados possam vir a contribuir para a construção de alternativas institucionais que venham a induzir a qualidade da educação superior.

% matched lemmas: arredondamento, avaliativo, coordenador, crescer, disjunto, dispensar, estimador, interseção, inverso, numérico, ordenar, parâmetro, pedagógico, percentagem, questionário, recorte, reiterar, reprodução, sequenciado, subconjunto, totalizar, vidar
\end{textsample}
